% Trust-Aware Explainable Federated Learning for Healthcare Risk Prediction
% Compiles with the official IEEEtran class when installed; otherwise falls
% back to the bundled ieee-fallback.sty, which reproduces an IEEE-like
% two-column conference layout.
\IfFileExists{IEEEtran.cls}
  {\documentclass[conference]{IEEEtran}\newcommand{\usingieeetran}{}}
  {\documentclass[10pt,twocolumn,letterpaper]{article}
   \usepackage{ieee-fallback}}

\usepackage{amsmath,amssymb}
\usepackage{graphicx}
\usepackage{booktabs}
\usepackage{url}
\usepackage{cite}
\graphicspath{{figures/}}

\begin{document}

\title{Trust-Aware Explainable Federated Learning for\\Personalized
Healthcare Risk Prediction}

% Six authors in a 3 x 2 grid. A nested tabular is used rather than
% \IEEEauthorblockN/\and so the block renders identically under IEEEtran and
% under the fallback layout.
\author{\small\setlength{\tabcolsep}{3pt}%
\begin{tabular}{@{}ccc@{}}
\begin{tabular}[t]{@{}c@{}}
  Dr H R Ravikumar\\
  \textit{Computer Science and Engineering}\\
  \textit{Sai Vidya Institute of Technology}\\
  Bengaluru, India\\
  ravikumar.hr@saividya.ac.in
\end{tabular} &
\begin{tabular}[t]{@{}c@{}}
  Asst. Prof. Lalitha K\\
  \textit{Computer Science and Engineering}\\
  \textit{Sai Vidya Institute of Technology}\\
  Bengaluru, India\\
  lalitha.k@saividya.ac.in
\end{tabular} &
\begin{tabular}[t]{@{}c@{}}
  Sahil Kumar\\
  \textit{Computer Science and Engineering}\\
  \textit{Sai Vidya Institute of Technology}\\
  Bengaluru, India\\
  sahil.kumar.23.cs@gmail.com
\end{tabular} \\ \noalign{\vskip 1.3em}
\begin{tabular}[t]{@{}c@{}}
  Mohit Khajuria\\
  \textit{Computer Science and Engineering}\\
  \textit{Sai Vidya Institute of Technology}\\
  Bengaluru, India\\
  iammohitkhajuria@gmail.com
\end{tabular} &
\begin{tabular}[t]{@{}c@{}}
  Manjunath K R\\
  \textit{Computer Science and Engineering}\\
  \textit{Sai Vidya Institute of Technology}\\
  Bengaluru, India\\
  manjunathkr14379@gmail.com
\end{tabular} &
\begin{tabular}[t]{@{}c@{}}
  Prem Kumar S\\
  \textit{Computer Science and Engineering}\\
  \textit{Sai Vidya Institute of Technology}\\
  Bengaluru, India\\
  premkumarsprem234@gmail.com
\end{tabular}
\end{tabular}}

\date{}
\maketitle

\begin{abstract}
Federated learning (FL) enables several hospitals to train a shared clinical
risk model without exchanging patient records, but standard aggregation such
as FedAvg assumes every participant is honest. A single institution that
submits poisoned updates can degrade the global model for all participants,
and conventional defences based only on update geometry become unreliable
when client data are strongly non-independent and identically distributed
(non-IID). This paper presents TA-EFL, a trust-aware and explainable
federated framework in which the server derives a per-client trust score from
the \emph{extrapolated attribution signature} of each submitted parameter
vector, evaluated on a shared public reference set, and aggregates updates in
proportion to $n_k T_k^{\gamma}$. Because the signature is recomputed
server-side from the weights a client actually submits, it cannot be spoofed
by reporting fabricated explanations. The framework is implemented on the
Flower FL framework with a logistic-regression risk model whose exact
closed-form SHAP values yield a personalized attribution for every
individual prediction rather than a single global feature ranking. Across five
tabular clinical prediction tasks and six experiments, trust-weighted
aggregation retains an area under the ROC curve (AUC) of $0.810 \pm 0.018$
with two of six clients malicious, where FedAvg falls to $0.769 \pm 0.067$;
an extrapolation factor of $\kappa=4$ raises the honest/attacker separation
margin to $+0.807$ from $+0.066$ at $\kappa=1$. A sensitivity sweep over the
trust-weighting parameter $\alpha$ shows that, \emph{in this experimental
configuration}, relying on explanation consistency alone ($\alpha=0$) gives
the largest separation margin, while AUC varies little across the sweep. All
datasets are public or synthetic and all hospitals are simulated; the system
is a research prototype and its outputs are risk estimates, not medical
diagnoses.
\end{abstract}

\begin{IEEEkeywords}
Federated learning, trust-aware aggregation, model poisoning, explainable
artificial intelligence, SHAP, healthcare risk prediction, non-IID data.
\end{IEEEkeywords}

\section{Introduction}

Clinical risk prediction models are most useful when they are trained on
data from many institutions, because a single hospital rarely observes a
patient population that is broad enough to support a model that generalises.
Patient records, however, are among the most tightly regulated categories of
personal data, and centralising them across institutional boundaries is often
legally and practically infeasible. Even where federation removes the need to
pool data, a recent scoping review finds that governance obligations do not
disappear but are redistributed across the participating
institutions \cite{eden2025governance}.

Federated learning (FL) addresses this tension by moving computation to the
data \cite{mcmahan2017,konecny2016}. Each participating hospital trains on
its own records and transmits only model parameters; a coordinating server
aggregates those parameters into a global model. FL has attracted sustained
interest in healthcare precisely because it offers multi-institution scale
without a central data repository \cite{rieke2020,sheller2020,dayan2021}. A recent systematic review of 612 studies confirms both the breadth of this activity and how little of it has reached deployment: only 5.2\% of the surveyed work reports a real-world application \cite{teo2024fl}.

Two problems limit the direct application of FL to clinical settings.

First, aggregation is trusting by construction. FedAvg \cite{mcmahan2017}
computes a sample-size weighted mean of client parameters and therefore
grants an adversarial participant influence in proportion to its claimed
dataset size. Label flipping, additive noise and scaled model poisoning are
all straightforward to implement by a participating institution
\cite{bhagoji2019,bagdasaryan2020,fang2020}, and recent surveys catalogue the resulting attack surface in detail \cite{sagar2023poisoning}. Byzantine-robust aggregation
rules such as Krum \cite{blanchard2017krum} and the coordinate-wise median
\cite{yin2018median} reduce this exposure by treating outlying updates as
suspicious, but they infer malice from the geometry of the update alone. In
healthcare, client distributions genuinely differ---a tertiary cardiac centre
and a rural clinic legitimately produce dissimilar updates---so
geometry-based outlier rejection risks penalising honest heterogeneity
\cite{hsieh2020noniid}.

Second, a model that influences clinical thinking must be interpretable. Post
hoc attribution methods, particularly SHAP \cite{lundberg2017shap}, grounded
in the Shapley value \cite{shapley1953}, are now standard for explaining
individual predictions, and the wider explainable AI literature argues that
high-stakes decisions require faithful explanations rather than plausible
narratives \cite{arrieta2020xai,doshivelez2017,rudin2019}. In federated
settings, explanations are usually treated purely as an output artefact
presented to clinicians after training has finished.

\textbf{Research gap.} These two concerns are typically addressed
separately: robustness work examines parameter geometry, and explainability
work examines the final model. The explanation signal is not used as evidence
about whether a client is behaving honestly. This is a missed opportunity,
because a poisoned model frequently attributes risk to the wrong features, or
with inverted sign, even when the norm and direction of its update look
unremarkable. A further difficulty is that a client cannot simply be asked to
report its own explanations, since a malicious participant would report
whatever explanation best conceals its behaviour.

\textbf{Contributions.} This paper makes the following contributions.
\begin{enumerate}
  \item A trust-aware federated aggregation scheme in which per-client trust
        is derived from \emph{explanation consistency} in addition to update
        geometry, combined through an exponential moving average and applied
        as aggregation weights $p_k \propto n_k T_k^{\gamma}$ with a hard
        exclusion threshold $\tau$.
  \item An \emph{extrapolated attribution signature}, computed entirely
        server-side by amplifying each client's update direction by a factor
        $\kappa$ before taking the signed attribution vector on a shared
        public reference set. Because it is derived from the submitted
        weights, it cannot be forged, and the amplification makes an
        inverted-attribution client detectable well before its update norm
        becomes anomalous.
  \item An implementation on the Flower framework
        \cite{beutel2020flower} with a desktop application that presents
        per-patient risk estimates together with both human-readable and
        technical SHAP explanations.
  \item An empirical study of six experiments spanning attack fraction,
        trust-component ablation, extrapolation sensitivity, heterogeneity
        false positives, cross-dataset generalisation, and sensitivity to
        the trust-weighting parameter, reported with per-seed variability
        and with negative results retained.
\end{enumerate}

\section{Related Work}

\subsection{Federated learning and aggregation}
FedAvg \cite{mcmahan2017} established the standard federated training loop:
broadcast the global model, train locally for several epochs, and average
the returned parameters weighted by local sample count. Surveys
\cite{kairouz2021,li2020spm} document its convergence behaviour and its
sensitivity to statistical heterogeneity; Hsieh \emph{et al.}
\cite{hsieh2020noniid} analyse the non-IID setting specifically. Flower
\cite{beutel2020flower} provides a framework-agnostic implementation with
pluggable server-side strategies, which is the mechanism this work extends.

\subsection{Trustworthiness in federated learning}
Surveys of trustworthy FL organise the field along axes of robustness,
privacy, fairness and explainability, and observe that these properties are
usually pursued in isolation \cite{tariq2023trustworthy}. The mechanism
proposed here couples two of them directly, using explainability as the
evidence on which robustness decisions are made.

\subsection{Poisoning and Byzantine-robust aggregation}
Model poisoning in FL has been demonstrated through adversarial-lens analysis
\cite{bhagoji2019}, backdoor insertion \cite{bagdasaryan2020} and attacks
crafted specifically against robust aggregators \cite{fang2020}. Defences
include Krum \cite{blanchard2017krum}, which selects the update closest to
its neighbours, and coordinate-wise median or trimmed mean
\cite{yin2018median}, which bound the influence of extreme coordinates. These
rules share an assumption: that a malicious update is geometrically atypical.
Fang \emph{et al.} \cite{fang2020} show that this assumption can be attacked
directly, and the taxonomy in \cite{sagar2023poisoning} organises these
attacks by adversary knowledge, target and stage of the training pipeline.
More recent defences acknowledge the same limitation from the opposite
direction: Molina-Coronado \cite{molina2024celtibero} reports that defences
which perform well on independent and identically distributed client data
often degrade once that assumption is relaxed, and proposes layered
aggregation in response.

\subsection{Explainable AI}
SHAP \cite{lundberg2017shap} unifies additive attribution methods under the
Shapley value \cite{shapley1953} and, for linear models, admits an exact
closed form. LIME \cite{ribeiro2016lime} provides local surrogate
explanations. Reviews \cite{arrieta2020xai,doshivelez2017} survey the field,
while Rudin \cite{rudin2019} argues that inherently interpretable models are
preferable for high-stakes decisions---an argument that motivates the choice
of a linear risk model in this work.

\subsection{Federated learning in healthcare}
Teo \emph{et al.} \cite{teo2024fl} systematically review clinical FL
applications and technical architectures. Rieke \emph{et al.}
\cite{rieke2020} set out the case for FL in digital health; Sheller \emph{et al.} \cite{sheller2020} demonstrate
multi-institutional training without data sharing; Dayan \emph{et al.}
\cite{dayan2021} report a large multi-national federated study for COVID-19
outcome prediction. Privacy analyses such as deep leakage from gradients
\cite{zhu2019dlg} show that transmitting parameters is not by itself
sufficient for privacy. On the engineering side, Fed-BioMed
\cite{cremonesi2023fedbiomed} argues that moving FL from simulation into
hospital practice is as much a software, networking and governance problem as
a modelling one---a gap that separates the present prototype from clinical
deployment.

\subsection{Gap addressed here}
A scoping review of the intersection between FL and explainable AI
\cite{lopezramos2024} finds that the two are usually combined so that
explanations are produced \emph{from} a federated model, rather than used
\emph{within} the federated protocol itself. Consistent with that finding,
existing trust or reputation mechanisms for FL
score clients using update statistics, validation loss on held-out data, or
historical behaviour, and treat explanations as a downstream artefact. This
work instead makes the attribution pattern of a submitted model a first-class
trust signal, computes it server-side so that it cannot be spoofed, and
amplifies it by extrapolation so that inverted attributions become
separable early in training.

\section{Proposed Methodology}

\subsection{System Architecture}
Fig.~\ref{fig:arch} shows the architecture. $K$ simulated hospitals hold
disjoint, non-IID shards of a tabular clinical dataset. Each round, the
Flower server broadcasts the global parameter vector, each client performs
local training and returns parameters, and the server computes trust scores,
aggregates, and evaluates the global model on a held-out test set. The
desktop application consumes the resulting model to produce per-patient risk
estimates and explanations.

\begin{figure}[t]
  \centering
  \includegraphics[width=\columnwidth]{architecture.png}
  \caption{TA-EFL architecture. Clients return only parameters; the server
  derives the extrapolated attribution signature, the trust score and the
  aggregation weights.}
  \label{fig:arch}
\end{figure}

\subsection{Federated Learning Framework}
The system is implemented against the Flower APIs
\cite{beutel2020flower}: each hospital is a \texttt{NumPyClient} exposed
through a \texttt{ClientApp}, and trust-weighted aggregation is a custom
server-side \texttt{Strategy} within a \texttt{ServerApp}. Two execution
backends are provided. The first is Flower's own simulation runtime. The
second advances the identical strategy object a single round per invocation,
which keeps experimental runs deterministic and allows the interface to
report progress while training is still in flight.

The local model is a binary logistic regression trained by mini-batch
stochastic gradient descent. For parameters $\theta = (\mathbf{w}, b)$,
\begin{equation}
  \hat{p}(\mathbf{x}) = \sigma(\mathbf{w}^{\top}\mathbf{x} + b),
  \qquad \sigma(z) = \frac{1}{1 + e^{-z}} .
\end{equation}
A linear model is chosen deliberately: it admits exact SHAP values in closed
form, so explanations introduce no approximation error of their own, in line
with the argument for interpretable models in high-stakes settings
\cite{rudin2019}.

The classical FedAvg baseline aggregates as
\begin{equation}
  \theta_g^{(t+1)} = \sum_{k=1}^{K} \frac{n_k}{\sum_{j} n_j}\,\theta_k^{(t)},
  \label{eq:fedavg}
\end{equation}
where $n_k$ is the number of local training samples at client $k$. The
coordinate-wise median baseline \cite{yin2018median} replaces this with an
elementwise median over $\{\theta_k^{(t)}\}$.

\subsection{Non-IID Client Partitioning}
Training indices are divided using a Dirichlet partition. For each class $c$,
proportions over clients are drawn as
\begin{equation}
  (q_{1,c},\dots,q_{K,c}) \sim \mathrm{Dir}(\alpha_{\mathrm{nIID}} \mathbf{1}_K),
\end{equation}
and class-$c$ indices are split accordingly, with resampling until every
client holds at least a minimum shard size. Smaller
$\alpha_{\mathrm{nIID}}$ yields more heterogeneous client populations. The
partition is a deterministic function of the dataset, client count,
$\alpha_{\mathrm{nIID}}$ and seed, so every simulated node reconstructs its
own shard without patient data being transferred between processes.

\subsection{Threat Model}
A subset $\mathcal{M} \subset \{1,\dots,K\}$ of hospitals is malicious. They
participate in the protocol correctly at the message level---returning
well-formed parameter vectors and truthful sample counts---but corrupt what
they return in one of three ways:
\begin{itemize}
  \item \textbf{Label flipping}: local labels are inverted, $y \mapsto 1-y$,
        before training.
  \item \textbf{Additive noise}: $\theta_k \mapsto \theta_k + \varepsilon$,
        $\varepsilon \sim \mathcal{N}(0, \sigma^{2} I)$ with $\sigma = 1.0$.
  \item \textbf{Scaled model poisoning}: $\theta_k \mapsto \theta_g +
        s(\theta_k - \theta_g)$ with $s = 8.0$.
\end{itemize}
The server is honest, observes only submitted parameters and sample counts,
and holds a small public reference set $X_{\mathrm{ref}}$ containing no
private records. Clients never transmit explanations, gradients on private
data, or raw features.

\subsection{Trust Computation}
Let $u_k = \theta_k - \theta_g$ be client $k$'s update and
$\tilde{u} = \mathrm{median}_k(u_k)$ the coordinate-wise median update. Two
similarity terms are computed. The update-similarity term is the cosine
similarity to the median update, attenuated by a magnitude penalty that
detects norm inflation:
\begin{equation}
  s^{\mathrm{upd}}_k = \max\!\left(0, \frac{u_k^{\top}\tilde{u}}
  {\lVert u_k\rVert\,\lVert \tilde{u}\rVert}\right)
  \cdot \min\!\left(1, \frac{\mathrm{median}_j \lVert u_j \rVert}
  {\lVert u_k \rVert}\right).
  \label{eq:supd}
\end{equation}
The explanation-consistency term compares attribution signatures (Section
\ref{sec:expltrust}):
\begin{equation}
  s^{\mathrm{exp}}_k = \max\!\left(0, \frac{\phi_k^{\top}\tilde{\phi}}
  {\lVert \phi_k\rVert\,\lVert \tilde{\phi}\rVert}\right),
  \qquad \tilde{\phi} = \mathrm{median}_k(\phi_k).
  \label{eq:sexp}
\end{equation}
The two are mixed by $\alpha$ and smoothed by an exponential moving average
with memory $\beta$:
\begin{equation}
  T_k(t) = \beta\,T_k(t-1) + (1-\beta)\left[\alpha s^{\mathrm{upd}}_k
  + (1-\alpha) s^{\mathrm{exp}}_k\right].
  \label{eq:trust}
\end{equation}
Thus $\alpha=1$ reduces trust to update geometry alone and $\alpha=0$ to
explanation consistency alone.

\subsection{Trust-Aware Aggregation}
Aggregation weights sharpen trust by an exponent $\gamma$ and exclude clients
below a threshold $\tau$:
\begin{equation}
  p_k = \frac{n_k \, \hat{T}_k^{\gamma}}{\sum_{j} n_j \hat{T}_j^{\gamma}},
  \qquad
  \hat{T}_k = \begin{cases} T_k, & T_k \geq \tau\\ 0, & T_k < \tau
  \end{cases}
  \label{eq:weights}
\end{equation}
with $\theta_g^{(t+1)} = \sum_k p_k \theta_k^{(t)}$. If every client falls
below $\tau$, the scheme falls back to uniform weighting rather than
producing an undefined update. Setting $T_k$ constant recovers
(\ref{eq:fedavg}).

\subsection{Explanation-Based Trust}
\label{sec:expltrust}
For a linear model, the exact SHAP value of feature $j$ at input
$\mathbf{x}$ in log-odds space is
\begin{equation}
  \varphi_j(\mathbf{x}) = w_j\left(x_j - \mathbb{E}[x_j]\right).
  \label{eq:shap}
\end{equation}
Aggregating over the shared reference set gives a per-client
\emph{signature}, an $L_2$-normalised signed vector,
\begin{equation}
  \phi_k = \frac{\mathbf{w}_k \odot \mathrm{sd}(X_{\mathrm{ref}})}
  {\lVert \mathbf{w}_k \odot \mathrm{sd}(X_{\mathrm{ref}})\rVert}.
  \label{eq:signature}
\end{equation}
The signature is signed rather than absolute, so a label-flipped model, whose
attributions point the wrong way, is anomalous even when its magnitude
profile is ordinary. Crucially, $\phi_k$ is computed by the server from the
submitted $\theta_k$, so a malicious client cannot report a favourable
explanation.

\subsection{$\kappa$ Extrapolation}
Early in training, individual updates are small and their signatures are
dominated by the shared global model, which compresses the difference
between honest and malicious clients. The server therefore extrapolates each
client's update before taking the signature:
\begin{equation}
  \theta_k^{\kappa} = \theta_g + \kappa\,(\theta_k - \theta_g),
  \qquad
  \phi_k = \mathrm{sig}\!\left(\theta_k^{\kappa}, X_{\mathrm{ref}}\right).
  \label{eq:kappa}
\end{equation}
This amplifies the direction in which a client is pushing the model without
altering the aggregation itself. Detection quality is quantified by the
\emph{detection margin}
\begin{equation}
  \Delta = \min_{k \notin \mathcal{M}} T_k - \max_{k \in \mathcal{M}} T_k,
  \label{eq:margin}
\end{equation}
which is positive only when every honest client is trusted more than every
attacker.

\subsection{SHAP Explainability}
Equation (\ref{eq:shap}) provides per-patient attributions at no
approximation cost. The application renders them twice: as a signed bar chart
over all features, and as a ranked plain-language summary that separates
factors increasing from factors decreasing the predicted risk, annotated with
the patient's own feature values in clinical units.

\subsection{Prediction and Graphical Interface}
A PyQt6 desktop application exposes federated training configuration, live
per-round trust diagnostics, model evaluation, per-patient explanation and
tabular reporting. Personalization here is at the level of the individual
patient record: the same global model is applied to every patient, but both
the risk estimate and the attribution breakdown are computed from that
patient's own feature vector, so two patients with the same predicted risk
may receive entirely different explanations. Predictions are presented as calibrated probabilities with
an explicit statement that they are risk estimates rather than diagnoses.

\section{Experimental Setup}

\subsection{Datasets}
Five binary-outcome tabular clinical tasks are used
(Table~\ref{tab:datasets}). The Breast Cancer Wisconsin (Diagnostic) data are
the real measurements distributed with scikit-learn \cite{pedregosa2011,
street1993wdbc}. For the remaining four tasks, the repository supports
loading the corresponding public CSV files
\cite{detrano1989,smith1988pima,dawber1951,dua2019uci}; when those files are
absent---as in the runs reported here---the framework generates a
calibrated synthetic cohort with realistic marginal distributions,
correlations and outcome prevalence. All reported results for those four
tasks are therefore obtained on synthetic data, and this is a limitation
rather than a claim of clinical validity (Section~\ref{sec:limits}).

\begin{table}[t]
  \caption{Datasets used in the experiments}
  \label{tab:datasets}
  \centering
  \small
  \setlength{\tabcolsep}{3.5pt}
  \begin{tabular}{lrrrl}
    \toprule
    Dataset & $n$ & Feat. & Pos. & Source \\
    \midrule
    Heart Disease       & 1500 & 13 & 52.6\% & Synthetic \\
    Pima Diabetes       & 1500 &  8 & 33.9\% & Synthetic \\
    Breast Cancer Wisc. &  569 & 10 & 37.3\% & Real (sklearn) \\
    Framingham CHD      & 2000 & 14 & 15.0\% & Synthetic \\
    Stroke Prediction   & 2500 & 10 &  9.8\% & Synthetic \\
    \bottomrule
  \end{tabular}
\end{table}

\subsection{Preprocessing and Missing Values}
Each dataset is split 75/25 into training and test partitions, stratified by
outcome. Features are standardised to zero mean and unit variance using
training-set statistics only, and the same statistics are applied to the test
partition and to any patient entered through the interface. For datasets
loaded from real CSV files, physiologically impossible zero entries in the
Pima data are treated as missing and imputed with the column median, and
categorical fields in the stroke data are integer-encoded; rows with residual
missing values are dropped. The synthetic generators do not produce missing
values.

\subsection{Federated Configuration}
Table~\ref{tab:config} lists the configuration used for all reported
experiments. The public reference set for attribution signatures contains
$\min(200, |X_{\mathrm{test}}|)$ samples drawn from the test pool; in a
deployment this would be a synthetic or public cohort rather than patient
records.

\begin{table}[t]
  \caption{Experimental configuration}
  \label{tab:config}
  \centering
  \small
  \begin{tabular}{ll}
    \toprule
    Parameter & Value \\
    \midrule
    Clients (hospitals) $K$        & 6 (experiments); 5 (interface default) \\
    Federated rounds               & 30 \\
    Local epochs / round           & 3 \\
    Learning rate, batch, $L_2$    & 0.15, 32, $10^{-3}$ \\
    Model                          & Logistic regression, mini-batch SGD \\
    Partition                      & Dirichlet, $\alpha_{\mathrm{nIID}}=0.5$ \\
    Train/test split               & 75\% / 25\%, stratified \\
    Trust weighting $\alpha$       & 0.0 (final, from E6) \\
    Extrapolation $\kappa$         & 4.0 \\
    EMA memory $\beta$             & 0.6 \\
    Weight exponent $\gamma$       & 2.0 \\
    Exclusion threshold $\tau$     & 0.25 \\
    Seeds                          & 42, 43, 44 \\
    Framework                      & Flower \cite{beutel2020flower} \\
    \bottomrule
  \end{tabular}
\end{table}

\subsection{Attacks, Baselines and Metrics}
Malicious clients are assigned by index so that the attacked configuration is
identical across strategies and seeds: one attacker is client~2 (label
flipping); two attackers add client~4 (noise); three attackers add client~1
(scaled poisoning). Trust-weighted aggregation is compared against FedAvg
(\ref{eq:fedavg}) and the coordinate-wise median \cite{yin2018median}.
Performance is reported as test-set ROC-AUC of the final global model,
averaged over three seeds with standard deviation; robustness of detection is
reported as the margin $\Delta$ in (\ref{eq:margin}).

\subsection{Selection of $\alpha$ and $\kappa$}
$\kappa$ was selected from the sweep in Section~\ref{sec:e3} and $\alpha$
from the sweep in Section~\ref{sec:e6}. Neither was tuned on a separate
criterion, and no other hyperparameter was adjusted while these sweeps were
performed.

\section{Results and Discussion}

Unless stated otherwise, results use the Heart Disease task with $K=6$,
30 rounds, and three seeds; means are reported with standard deviation.
Sections~\ref{sec:e1}, \ref{sec:e3} and \ref{sec:e5} use the final
$\alpha=0$ configuration.

\subsection{E1: Malicious-Client Fraction}
\label{sec:e1}
Table~\ref{tab:e1} and Fig.~\ref{fig:e1} report final AUC as the fraction of
malicious clients increases. With no attacker, all three aggregation rules
are equivalent, as expected: trust weighting does not help when there is
nothing to defend against. The strategies separate once attackers are
present. At two of six malicious clients, FedAvg falls to
$0.769 \pm 0.067$ while trust weighting holds $0.810 \pm 0.018$; the
increased standard deviation of FedAvg indicates that the damage depends
strongly on the particular partition. The coordinate-wise median remains
competitive throughout, which is consistent with its design: it is a robust
estimator, and the attacks used here are not tailored to defeat it
\cite{fang2020}.

\begin{table}[t]
  \caption{E1: final AUC vs.\ malicious fraction (mean $\pm$ sd, 3 seeds)}
  \label{tab:e1}
  \centering
  \small
  \begin{tabular}{lccc}
    \toprule
    Malicious & Trust-Aware & FedAvg & Median \\
    \midrule
    0/6 & $0.808 \pm 0.015$ & $0.808 \pm 0.014$ & $0.807 \pm 0.018$ \\
    1/6 & $\mathbf{0.807 \pm 0.016}$ & $0.806 \pm 0.015$ & $0.796 \pm 0.019$ \\
    2/6 & $\mathbf{0.810 \pm 0.018}$ & $0.769 \pm 0.067$ & $0.802 \pm 0.013$ \\
    3/6 & $\mathbf{0.806 \pm 0.020}$ & $0.775 \pm 0.059$ & $0.789 \pm 0.036$ \\
    \bottomrule
  \end{tabular}
\end{table}

\begin{figure}[t]
  \centering
  \includegraphics[width=\columnwidth]{e1_robustness.png}
  \caption{E1: robustness to increasing attack fraction, $\alpha=0$.}
  \label{fig:e1}
\end{figure}

\subsection{E2: Trust-Component Ablation}
\label{sec:e2}
Table~\ref{tab:e2} isolates the two trust signals at two of six malicious
clients. Using update similarity alone ($\alpha=1$) yields a \emph{negative}
margin, $-0.095 \pm 0.127$: at least one attacker is trusted more than at
least one honest client, so the ordering is not usable for exclusion. Adding
explanation consistency reverses this, and using explanation consistency
alone gives the largest margin, $+0.513 \pm 0.296$. AUC differences among the
three settings are small relative to seed variance. This experiment sets
$\alpha$ explicitly and is therefore unaffected by the choice of default.

\begin{table}[t]
  \caption{E2: trust-component ablation, 2/6 malicious (3 seeds)}
  \label{tab:e2}
  \centering
  \small
  \begin{tabular}{lcc}
    \toprule
    Trust signal & AUC & Detection margin $\Delta$ \\
    \midrule
    Update only ($\alpha=1$)      & $0.800 \pm 0.017$ & $-0.095 \pm 0.127$ \\
    Both ($\alpha=0.5$)           & $0.805 \pm 0.013$ & $+0.243 \pm 0.111$ \\
    Explanation only ($\alpha=0$) & $\mathbf{0.810 \pm 0.018}$ &
                                    $\mathbf{+0.513 \pm 0.296}$ \\
    \bottomrule
  \end{tabular}
\end{table}

\begin{figure}[t]
  \centering
  \includegraphics[width=\columnwidth]{e2_ablation.png}
  \caption{E2: contribution of the two trust signals.}
  \label{fig:e2}
\end{figure}

\subsection{E3: Extrapolation Sensitivity}
\label{sec:e3}
Table~\ref{tab:e3} and Fig.~\ref{fig:e3} sweep $\kappa$ with a single
label-flipping client. Without meaningful extrapolation ($\kappa=1$) the
margin is $+0.066 \pm 0.048$, close to the point at which honest and
malicious clients are indistinguishable. The margin peaks at $\kappa=4$
($+0.807 \pm 0.080$) and declines for larger values, where extrapolation
begins to amplify noise as well as signal. This supports the central claim
that amplifying the update direction, rather than inspecting the raw update,
is what exposes inverted attribution. Note that $\kappa$ affects only trust
scoring; the aggregation in (\ref{eq:weights}) uses unmodified parameters.

\begin{table}[t]
  \caption{E3: detection margin vs.\ $\kappa$, one label flipper (3 seeds)}
  \label{tab:e3}
  \centering
  \small
  \begin{tabular}{lc}
    \toprule
    Extrapolation $\kappa$ & Detection margin $\Delta$ \\
    \midrule
    1  & $+0.066 \pm 0.048$ \\
    2  & $+0.674 \pm 0.399$ \\
    4  & $\mathbf{+0.807 \pm 0.080}$ \\
    8  & $+0.495 \pm 0.190$ \\
    16 & $+0.139 \pm 0.073$ \\
    \bottomrule
  \end{tabular}
\end{table}

\begin{figure}[t]
  \centering
  \includegraphics[width=\columnwidth]{e3_kappa.png}
  \caption{E3: detection margin as a function of $\kappa$, $\alpha=0$.}
  \label{fig:e3}
\end{figure}

\subsection{E4: Heterogeneity False Positives}
\label{sec:e4}
A robustness mechanism is only usable if it does not penalise honest
institutions whose patient populations differ. E4 runs with \emph{no}
attackers and varies $\alpha_{\mathrm{nIID}}$, recording the minimum trust
assigned to any honest client (Table~\ref{tab:e4}). Across the range tested,
the minimum honest trust remains between $0.384$ and $0.456$, above the
exclusion threshold $\tau = 0.25$, so no honest client is excluded even at
$\alpha_{\mathrm{nIID}} = 0.1$. These values were measured under the earlier
$\alpha = 0.5$ configuration and have not been re-measured at $\alpha = 0$;
they should therefore be read as evidence for the mixed-signal setting.

\begin{table}[t]
  \caption{E4: minimum honest-client trust, no attackers (3 seeds)}
  \label{tab:e4}
  \centering
  \small
  \begin{tabular}{lc}
    \toprule
    Dirichlet $\alpha_{\mathrm{nIID}}$ & Min.\ honest trust \\
    \midrule
    0.1 & $0.432 \pm 0.023$ \\
    0.5 & $0.411 \pm 0.025$ \\
    1.0 & $0.384 \pm 0.065$ \\
    5.0 & $0.456 \pm 0.013$ \\
    \bottomrule
  \end{tabular}
\end{table}

\subsection{E5: Cross-Dataset Generalisation}
\label{sec:e5}
Table~\ref{tab:e5} repeats the two-of-six attack configuration on all five
tasks. Trust weighting attains the highest or joint-highest AUC on every
dataset. The largest gap appears on Breast Cancer Wisconsin, where FedAvg
collapses to $0.657$: the label-flipping client inverts a model that is
otherwise highly separable, and sample-size weighting propagates that
inversion. Detection margins are positive on all five tasks, though smallest
on Stroke Prediction ($+0.199$), the most imbalanced task, where honest
clients disagree most and the explanation signal is correspondingly weaker.

\begin{table}[t]
  \caption{E5: cross-dataset AUC and detection margin, 2/6 malicious}
  \label{tab:e5}
  \centering
  \small
  \begin{tabular}{lcccc}
    \toprule
    Dataset & Trust & FedAvg & Median & $\Delta$ \\
    \midrule
    Heart Disease       & $\mathbf{0.810}$ & 0.769 & 0.802 & $+0.513$ \\
    Pima Diabetes       & $\mathbf{0.895}$ & 0.874 & 0.891 & $+0.688$ \\
    Breast Cancer Wisc. & $\mathbf{0.977}$ & 0.657 & $\mathbf{0.977}$ &
                          $+0.707$ \\
    Framingham CHD      & $\mathbf{0.796}$ & 0.720 & 0.785 & $+0.484$ \\
    Stroke Prediction   & $\mathbf{0.909}$ & 0.828 & 0.891 & $+0.199$ \\
    \bottomrule
  \end{tabular}
\end{table}

\begin{figure}[t]
  \centering
  \includegraphics[width=\columnwidth]{e5_datasets.png}
  \caption{E5: cross-dataset robustness and honest/attacker separation.}
  \label{fig:e5}
\end{figure}

\subsection{E6: Sensitivity to the Trust Weighting $\alpha$}
\label{sec:e6}
Table~\ref{tab:e6} and Fig.~\ref{fig:e6} sweep $\alpha$ over
$\{0, 0.25, 0.5, 0.75, 1\}$ under the E2 protocol: same dataset, attack
configuration, partitions, seeds, $\kappa$ and evaluation, with $\alpha$ the
only variable.

Three observations follow. First, the detection margin decreases
monotonically as the update-similarity contribution grows, from
$+0.514 \pm 0.296$ at $\alpha=0$ to $-0.095 \pm 0.127$ at $\alpha=1$, where
separation is lost entirely. Second, AUC varies comparatively little across
the sweep---from $0.800$ to $0.810$, a range of the same order as the
per-seed standard deviation---so AUC alone does not discriminate strongly
between settings. Third, $\alpha=0$ attains both the largest margin and the
highest mean AUC in this experimental environment.

$\alpha=0$ was therefore adopted for the final configuration \emph{primarily
on the basis of malicious-client separation}. This should not be read as a
claim that $\alpha=0$ is universally optimal. The result is specific to these
datasets, this attack set, this model family and this partitioning; the
update-similarity term retains a clear purpose against attacks that preserve
attribution structure while inflating magnitude, and a different attack
mixture could favour a non-zero $\alpha$. What the sweep does establish is
that, here, explanation consistency carries the discriminative signal and
adding update similarity dilutes it.

\begin{table}[t]
  \caption{E6: sensitivity to trust weighting $\alpha$, 2/6 malicious
  (3 seeds)}
  \label{tab:e6}
  \centering
  \small
  \begin{tabular}{lcc}
    \toprule
    $\alpha$ & AUC & Detection margin $\Delta$ \\
    \midrule
    0.00 & $\mathbf{0.810 \pm 0.018}$ & $\mathbf{+0.514 \pm 0.296}$ \\
    0.25 & $0.810 \pm 0.018$ & $+0.413 \pm 0.187$ \\
    0.50 & $0.805 \pm 0.013$ & $+0.243 \pm 0.111$ \\
    0.75 & $0.805 \pm 0.014$ & $+0.061 \pm 0.111$ \\
    1.00 & $0.800 \pm 0.017$ & $-0.095 \pm 0.127$ \\
    \bottomrule
  \end{tabular}
\end{table}

\begin{figure}[t]
  \centering
  \includegraphics[width=\columnwidth]{e6_alpha.png}
  \caption{E6: $\alpha$ against AUC (left) and detection margin (right).}
  \label{fig:e6}
\end{figure}

\section{Explainability and Graphical Interface}

The application presents the prediction before any explanation. For a
selected patient it reports a risk category and the estimated probability,
followed by a ranked breakdown of the factors that increased and decreased
that estimate, each annotated with the patient's own value in clinical units
and a qualitative strength. A plain-language summary names the strongest
contributor on each side. Below this, the technical view shows the signed
SHAP attribution over all features (Fig.~\ref{fig:shap}), with the axis
labelled as impact on the prediction rather than as a clinical effect.

\begin{figure}[t]
  \centering
  \includegraphics[width=\columnwidth]{shap_example.png}
  \caption{Per-patient attribution from (\ref{eq:shap}) for one test patient.
  Positive values increase the predicted risk. Illustrative example produced
  by the project's explainability module on the synthetic Heart Disease
  cohort.}
  \label{fig:shap}
\end{figure}

Categorical inputs are presented by name rather than by code---sex as
female/male, chest pain type, resting ECG, ST slope and thalassemia by their
category labels---while the numeric code required by the model is retained
internally, so the vector reaching the model is unchanged. The same interface
can explain either a held-out test patient or the patient just entered on the
prediction screen; in the latter case the explanation and the prediction
refer to the same input by construction.

The reporting screen presents model performance, attack robustness,
per-client trust decomposition, the E1--E6 summary and cross-dataset results
as tables, reading from the stored experiment outputs rather than recomputing
them, and supports CSV export.

Two framing decisions are deliberate. Attributions are described as
contributions of features to the model's output, not as causal or clinical
statements, following the caution urged in the explainability literature
\cite{arrieta2020xai,rudin2019}. Every prediction view carries an explicit
notice that the output is an AI-generated risk estimate and not a medical
diagnosis.

\section{Limitations}
\label{sec:limits}

The following limitations bound the claims made here.

\emph{Simulated federation.} All hospitals are simulated processes on one
machine. Real deployments face stragglers, dropout, heterogeneous compute and
communication failures, none of which are modelled. They also face governance
questions---consent, liability, auditability and data stewardship across
institutions---that a technical mechanism does not resolve
\cite{eden2025governance}.

\emph{Data provenance.} Only the Breast Cancer Wisconsin task uses real
measurements in the reported runs. The other four use calibrated synthetic
cohorts, which reproduce plausible marginals and correlations but not the
messiness, coding drift or site-specific bias of hospital records. Results on
those four tasks characterise the method's behaviour, not clinical
performance.

\emph{Dataset scale.} Cohorts of 569--2500 records split across six clients
leave some shards small, which inflates seed-to-seed variance; several
reported standard deviations are large relative to the differences discussed.

\emph{Attack scope.} Three attack types are evaluated. Adaptive adversaries
that optimise against the trust rule itself---for example, by matching the
cohort attribution signature while poisoning slowly---are not evaluated, and
\cite{fang2020} shows such attacks are practical against robust aggregators.

\emph{Explanation fidelity.} Exact SHAP values are available only because
the model is linear. The signature in (\ref{eq:signature}) summarises a model
by its weight profile on a reference set, so two models with different
behaviour but similar weight direction would be scored alike. SHAP values
also describe association within the model, not causation.

\emph{No clinical validation.} The system has not been evaluated
prospectively, on real patient cohorts, or in any care setting; no clinical
utility, safety or equity claim is made. Its outputs are risk estimates from
a research prototype and must not be used for diagnosis or treatment
decisions.

\emph{Environment sensitivity.} Absolute values depend on the numerical
environment: the split implementation and metric library affect partitioning
and hence results. All figures reported here come from a single consistent
environment and configuration.

\section{Conclusion and Future Work}

This paper presented TA-EFL, a federated learning framework in which
explanation consistency is used as a trust signal rather than only as an
output. By recomputing an extrapolated attribution signature server-side from
each submitted parameter vector, the method detects clients whose models
attribute risk in the wrong direction without requiring clients to report
anything beyond parameters. Across six experiments, trust-weighted
aggregation preserved AUC under attack where FedAvg degraded, extrapolation
at $\kappa=4$ produced the strongest honest/attacker separation, and no
honest client was excluded under the heterogeneity levels tested. A
sensitivity sweep led to $\alpha=0$ for this configuration, selected on
separation rather than on accuracy, and reported as configuration-specific
rather than universal.

Future work follows directly from the limitations. The most important step is
evaluation on real multi-institutional data with a clinical partner. Beyond
that: adaptive attacks that target the attribution signature itself; non-linear
models, where exact SHAP is unavailable and approximate attributions would
introduce their own error; formal privacy guarantees for the reference set and
the signatures, given known gradient-leakage risks \cite{zhu2019dlg};
re-measurement of E4 under the final configuration; and calibration analysis,
since a risk score intended to inform clinicians must be well calibrated and
not merely well ranked.

\section*{Reproducibility}

All experiments are reproducible from the accompanying repository with fixed
seeds $\{42,43,44\}$. Experiments E1, E3 and E5 in the final configuration
are produced by \texttt{run\_study.py --only e1,e3,e5 --label final\_alpha0};
E2 and E4 by \texttt{run\_study.py}; E6 by
\texttt{run\_study.py --only e6}. Per-seed values are stored as CSV files
alongside the figures reproduced in this paper.

\bibliographystyle{ieeetr}
\bibliography{refs}

\end{document}
